\documentclass[11pt,a4paper]{article}

\usepackage[utf8]{inputenc}
\usepackage[T1]{fontenc}
\usepackage{XCharter}
\usepackage[brazil]{babel}
\usepackage[margin=2cm]{geometry}
\usepackage{amsmath, amssymb}
\usepackage[xcharter,bigdelims,vvarbb]{newtxmath}
% newtxmath's operators font requests the fontspec family `XCharter(0)' in OT1;
% without these substitutions math digits and operator names silently fall back
% to Computer Modern (LaTeX Font Warning: Font shape `OT1/XCharter(0)/m/n' undefined).
\DeclareFontFamily{OT1}{XCharter(0)}{}
\DeclareFontShape{OT1}{XCharter(0)}{m}{n}{<->ssub * XCharter-TLF/m/n}{}
\DeclareFontShape{OT1}{XCharter(0)}{m}{it}{<->ssub * XCharter-TLF/m/it}{}
\DeclareFontShape{OT1}{XCharter(0)}{b}{n}{<->ssub * XCharter-TLF/b/n}{}
\DeclareFontShape{OT1}{XCharter(0)}{bx}{n}{<->ssub * XCharter-TLF/b/n}{}
\usepackage{enumerate}
\usepackage{xcolor}

\pagestyle{empty}
\setlength{\parindent}{0pt}
\setlength{\parskip}{0.5em}

\newcommand{\thetav}{\boldsymbol{\theta}}

\begin{document}

%% =============================================================
%%  Header
%% =============================================================
{\Large\bfseries Aprendizado Profundo \textemdash{} Gabarito} \hfill \textbf{Prova G2, Unidades 4 e 5}\\
Prof. Lucas S. Kupssinskü

\rule{\linewidth}{0.8pt}

%% =============================================================
%%  Objective questions
%% =============================================================
\section*{Questões Objetivas}

\textbf{Quadro de respostas:}
\begin{center}
\begin{tabular}{|c|c|c|c|c|c|c|c|}
\hline
1 & 2 & 3 & 4 & 5 & 6 & 7 & 8 \\\hline
e & d & c & g & d & a & h & f \\\hline
\end{tabular}
\end{center}

\color{blue}

\subsection*{Questão 1 \textemdash{} Resposta: (e)}
Cada matriz adaptada recebe $\mathbf{B} \in \mathbb{R}^{1024 \times 8}$ e $\mathbf{A} \in \mathbb{R}^{8 \times 1024}$:
\[
r(d + k) = 8 \times (1024 + 1024) = 16\,384 \text{ parâmetros por matriz adaptada}.
\]
São $2$ projeções ($\mathbf{W}_Q$ e $\mathbf{W}_V$) em cada uma das $12$ camadas:
\[
16\,384 \times 2 \times 12 = 393\,216.
\]
\textbf{(a)} $16\,384$ conta uma única projeção de uma única camada.
\textbf{(b)} $32\,768$ conta as duas projeções de uma única camada.
\textbf{(c)} $98\,304$ conta apenas $B$ (ou apenas $A$) de uma projeção nas $12$ camadas.
\textbf{(d)} $196\,608$ esquece uma das duas projeções ($16\,384 \times 12$).
\textbf{(f)} $589\,824$ aplica a três projeções ($\mathbf{Q}$, $\mathbf{K}$ e $\mathbf{V}$), não às duas do enunciado.
\textbf{(g)} $786\,432$ aplica às quatro projeções ($\mathbf{Q}$, $\mathbf{K}$, $\mathbf{V}$ e $\mathbf{O}$).
\textbf{(h)} $25\,165\,824$ treina o $\Delta\mathbf{W}$ denso $1024 \times 1024$ em vez dos fatores \textit{low rank}, exatamente o que o LoRA evita.

\subsection*{Questão 2 \textemdash{} Resposta: (d)}
\textbf{I é correta:} no SFT a perda $\mathcal{L}_{SFT} = -\sum_{t \in R} \log P_\theta(x_t \mid x_{<t})$ soma apenas sobre os \textit{tokens} da resposta $R$; o \textit{prompt} entra como contexto, mas não gera gradiente.

\textbf{II é correta:} o DPO deriva a recompensa implícita da própria política e otimiza uma perda de classificação sobre pares $(y_w, y_l)$, eliminando o \textit{reward model} explícito e a amostragem \textit{online} do loop de PPO.

\textbf{III é incorreta:} o termo $\mathrm{KL}(\pi_\theta \| \pi_{\text{ref}})$ não acelera a convergência; ele \emph{restringe} a política a permanecer próxima da referência, evitando que o modelo degenere ou explore falhas do \textit{reward model} (\textit{reward hacking}).

Logo, apenas I e II estão corretas.\\
\textbf{(a)/(b)} excluem uma asserção correta.
\textbf{(c)/(e)/(f)/(g)} incluem a asserção III, que inverte a função do termo de KL.
\textbf{(h)} contradiz I e II.

\subsection*{Questão 3 \textemdash{} Resposta: (c)}
O quadro mostrado é sintoma do \textbf{esquecimento catastrófico}: a acurácia na tarefa alvo sobe enquanto capacidades adquiridas no pré-treino se degradam, porque o \textit{fine-tuning} completo com taxa alta e muitas épocas sobrescreve os pesos que codificavam o conhecimento geral. Mitigações padrão: reduzir taxa/épocas, misturar dados gerais ao corpus de ajuste (\textit{replay}) ou usar PEFT (LoRA), que congela $\mathbf{W}_0$ e limita a atualização a poucos parâmetros.\\
\textbf{(a)} a curva jurídica é medida em dados de teste não vistos no treino e sobe de forma consistente, ou seja, o modelo \emph{generaliza} no domínio alvo; \textit{overfitting} não explica esse padrão nem a queda no \textit{benchmark} geral, e mais épocas agravariam o quadro.
\textbf{(b)} vazamento de dados inflaria a curva jurídica, mas não derrubaria o \textit{benchmark} geral.
\textbf{(d)} \textit{reward hacking} pertence ao contexto de RLHF com \textit{reward model}, ausente aqui.
\textbf{(e)} um \textit{tokenizer} ruim para o domínio prejudicaria a curva jurídica, não a geral; retreinar o \textit{tokenizer} invalidaria os \textit{embeddings} do pré-treino.
\textbf{(f)} janela de contexto excedida causa truncamento nas entradas jurídicas, sem efeito sobre o \textit{benchmark} geral.
\textbf{(g)} explosão de gradiente com \texttt{NaN} destruiria as duas curvas; o gráfico mostra melhora consistente na tarefa alvo.
\textbf{(h)} leis de escala tratam do pré-treino; adicionar dados jurídicos não recupera o conhecimento geral perdido.

\subsection*{Questão 4 \textemdash{} Resposta: (g)}
Associação correta: \textbf{1-(ii)} o \textit{PE} sinusoidal é uma função fixa da posição (senos e cossenos em frequências $10000^{-2k/d}$), somada à entrada, sem parâmetros; \textbf{2-(i)} o \textit{PE} absoluto aprendido é uma tabela $\mathbf{E} \in \mathbb{R}^{L_{\text{max}} \times d}$ treinável, que não extrapola além de $L_{\text{max}}$; \textbf{3-(iv)} o RoPE aplica rotações $\mathbf{R}_\theta(i)$ a $\mathbf{q}$ e $\mathbf{k}$ dentro da atenção, e o \textit{score} resultante depende só de $j - i$; \textbf{4-(iii)} o ALiBi subtrai $m_h\,|i - j|$ do \textit{score}, com inclinação $m_h$ própria por cabeça.

Erros agrupados por troca: \textbf{(a)} e \textbf{(d)} trocam RoPE com ALiBi, confundindo rotação interna com penalidade linear; \textbf{(b)} e \textbf{(d)} trocam sinusoidal com aprendido, ignorando que só o segundo tem parâmetros e limite $L_{\text{max}}$; \textbf{(c)} atribui a tabela treinável ao RoPE e a rotação ao \textit{PE} aprendido; \textbf{(e)} troca sinusoidal com RoPE; \textbf{(f)} atribui a penalidade por distância ao \textit{PE} aprendido; \textbf{(h)} troca sinusoidal com ALiBi, atribuindo ao primeiro um mecanismo interno à atenção.

\subsection*{Questão 5 \textemdash{} Resposta: (d)}
Pela regra da cadeia, com $\sigma'(d) = \sigma(d)(1 - \sigma(d))$:
\[
\frac{\partial L}{\partial d} = \frac{-\sigma'(d)}{1 - \sigma(d)} = -\sigma(d) = -0{,}05, \qquad
\frac{\partial L_{\text{ns}}}{\partial d} = \frac{-\sigma'(d)}{\sigma(d)} = -(1 - \sigma(d)) = -0{,}95.
\]
Logo $\left|\partial L / \partial d\right| = 0{,}05$ e $\left|\partial L_{\text{ns}} / \partial d\right| = 0{,}95$: quando o discriminador vence ($\sigma \to 0$), a \textit{loss} original satura (gradiente $\propto \sigma$) e o gerador para de aprender, enquanto a \textit{non-saturating} mantém gradiente $\propto 1 - \sigma$, próximo de $1$. É exatamente a motivação da troca de \textit{loss}.\\
\textbf{(a)} acerta os valores mas inverte a leitura: sinal maior vem do gradiente de maior módulo, o da \textit{non-saturating}.
\textbf{(b)/(e)} trocam as derivadas entre as duas \textit{losses}.
\textbf{(c)} aplica a derivada da \textit{loss} original às duas.
\textbf{(f)} esquece o denominador da regra da cadeia e fica com $\sigma'(d) = 0{,}0475$ para ambas.
\textbf{(g)} aplica a derivada da \textit{non-saturating} às duas.
\textbf{(h)} trata $\sigma(d) = 0{,}05$ como exatamente zero; ambos os gradientes são finitos e bem definidos nesse ponto.

\subsection*{Questão 6 \textemdash{} Resposta: (a)}
\textbf{I é verdadeira:} no treinamento com \textit{classifier-free guidance}, descarta-se a condição $c$ (substituída por $\varnothing$) em uma fração pequena dos exemplos, tipicamente $p \approx 0{,}1$.

\textbf{II é verdadeira e justifica I:} o descondicionamento esporádico existe exatamente para que a \emph{mesma} rede aprenda as duas predições, condicional e incondicional, sem custo de um segundo modelo; na amostragem elas são combinadas por $\tilde{\mathbf{g}}_t = (1+w)\,\mathbf{g}_t[\mathbf{z}_t, c] - w\,\mathbf{g}_t[\mathbf{z}_t, \varnothing]$ com $w \geq 0$, extrapolando na direção da condição.\\
\textbf{(b)} a relação causal existe: II descreve a finalidade do procedimento descrito em I.
\textbf{(c)} inverte a direção: o descondicionamento no treino (I) é o meio; a combinação na amostragem (II) é o fim.
\textbf{(d)/(f)/(h)} negam a asserção II, que reproduz a formulação vista em aula.
\textbf{(e)} nega a asserção I, que descreve o protocolo de treinamento padrão do \textit{classifier-free guidance}.
\textbf{(g)} com $w = 0$ a amostragem usa só a predição condicional, mas o treino descondicionado independe do $w$ escolhido depois; a condição inventada não existe.
\textbf{(h)} difusão latente muda o espaço em que a difusão ocorre, não o protocolo de guiamento.

\subsection*{Questão 7 \textemdash{} Resposta: (h)}
Ordem válida: \textbf{I, II, III, IV, V}. As dependências: o kernel (III) precisa de $\mathbf{x}$ (I) e de $t, \boldsymbol{\epsilon}$ (II); a perda (IV) precisa de $\mathbf{z}_t$ (III) e de $\boldsymbol{\epsilon}$ para comparar; a atualização (V) precisa do gradiente da perda (IV).\\
\textbf{(a)} computa o kernel antes de amostrar $t$ e $\boldsymbol{\epsilon}$, dos quais ele depende.
\textbf{(b)} computa o kernel antes de amostrar a instância $\mathbf{x}$.
\textbf{(c)} computa a perda antes de construir $\mathbf{z}_t$.
\textbf{(d)} inicia pelo kernel, sem $\mathbf{x}$, $t$ ou $\boldsymbol{\epsilon}$ amostrados.
\textbf{(e)} atualiza os parâmetros antes de computar a perda.
\textbf{(f)} atualiza os parâmetros antes de qualquer perda; não há gradiente acumulado de iterações anteriores no algoritmo do DDPM.
\textbf{(g)} computa a perda antes do kernel, com $\mathbf{z}_t$ ainda inexistente.

\subsection*{Questão 8 \textemdash{} Resposta: (f)}
Pela mudança de variável, $p_Y(y) = p_X\!\left(\tfrac{y - 1}{3}\right) \cdot \left|\tfrac{dX}{dY}\right| = p_X\!\left(\tfrac{y - 1}{3}\right) \cdot \tfrac{1}{3}$. Em $y = 1$:
\[
p_Y(1) = p_X(0) \cdot \tfrac{1}{3} \approx \frac{0{,}399}{3} \approx 0{,}133.
\]
\textbf{(a)} $1{,}197$ multiplica por $3$ em vez de dividir; a transformação \emph{estica} o espaço, logo a densidade deve \emph{diminuir}.
\textbf{(b)} $0{,}798$ multiplica por $2$: coeficiente errado e direção errada da correção de volume.
\textbf{(c)} $0{,}399$ ignora o Jacobiano por completo.
\textbf{(d)} $0{,}242$ é $p_X(1)$: avalia no ponto errado e ignora o Jacobiano.
\textbf{(e)} $0{,}199$ divide por $2$, reproduzindo o coeficiente do exemplo visto em aula em vez do da questão.
\textbf{(g)} $0{,}081$ avalia a densidade base em $y = 1$ em vez de em $x = (y-1)/3 = 0$ ($p_X(1)/3$).
\textbf{(h)} $0{,}044$ divide por $9$, elevando o Jacobiano ao quadrado.

\color{black}
\rule{\linewidth}{0.4pt}

%% =============================================================
%%  Dissertative questions
%% =============================================================
\section*{Questões Dissertativas}

\color{blue}

\subsection*{Questão 9 \textemdash{} \textit{Reward model} e DPO}

\textbf{(a)} Modelo de Bradley--Terry:
\[
P(y_w \succ y_l) = \sigma\!\left(r_\phi(x, y_w) - r_\phi(x, y_l)\right) = \sigma(2{,}2 - 0{,}7) = \sigma(1{,}5) = \frac{1}{1 + e^{-1{,}5}} \approx 0{,}8176.
\]
\[
\mathcal{L}_{RM} = -\log \sigma(1{,}5) \approx -\ln(0{,}8176) \approx 0{,}2014.
\]

\textbf{(b)} DPO:
\[
\beta\left[1{,}5 - (-1{,}0)\right] = 0{,}2 \cdot 2{,}5 = 0{,}5, \qquad
\mathcal{L}_{DPO} = -\log \sigma(0{,}5) = -\ln(0{,}6225) \approx 0{,}4741.
\]
\textit{Explicação esperada:} o DPO dispensa o treinamento de um \textit{reward model} explícito e o loop de \textit{reinforcement learning} (PPO) com amostragem \textit{online}; a recompensa fica implícita nos log-\textit{ratios} entre a política e a referência, e o alinhamento vira uma perda supervisionada de classificação sobre os pares de preferência.

\subsection*{Questão 10 \textemdash{} Difusão: kernel e passo reverso}

\textbf{(a)} \textit{Schedule} acumulado:
\[
\alpha_1 = 1 - \beta_1 = 0{,}6, \qquad \alpha_2 = (1 - \beta_1)(1 - \beta_2) = 0{,}6 \cdot 0{,}75 = 0{,}45.
\]
Kernel de difusão:
\[
z_2 = \sqrt{\alpha_2}\,x + \sqrt{1 - \alpha_2}\,\epsilon = \sqrt{0{,}45} \cdot 2 + \sqrt{0{,}55} \cdot 0{,}5 \approx 1{,}3416 + 0{,}3708 = 1{,}7125.
\]

\textbf{(b)} Amostragem reversa com $\sigma_2 = 0$:
\[
\hat{z}_1 = \frac{1}{\sqrt{1 - \beta_2}}\left(z_2 - \frac{\beta_2}{\sqrt{1 - \alpha_2}}\,g_2[z_2, \thetav_2]\right)
= \frac{1}{\sqrt{0{,}75}}\left(1{,}7125 - \frac{0{,}25}{0{,}7416} \cdot 0{,}6\right)
\]
\[
= 1{,}1547 \cdot (1{,}7125 - 0{,}2023) = 1{,}1547 \cdot 1{,}5102 \approx 1{,}7438.
\]
\textit{Explicação esperada:} o termo $\sigma_t\,\epsilon$ injeta estocasticidade no processo reverso: cada $z_{t-1}$ é \emph{amostrado} de uma gaussiana centrada na predição da rede, o que permite gerar amostras diversas a partir do mesmo $z_T$. Com $\sigma_t = 0$ a trajetória reversa torna-se determinística.

\color{black}

\end{document}
